\subsection{Логические операции}

\textbf{Мотивирующий вопрос:} Как дрон понимает, что <<батарея разряжена>> И <<высота больше нуля>>?

Логические операторы позволяют комбинировать условия. Результатом всегда является \texttt{True} (Истина) или \texttt{False} (Ложь).

\subsubsection*{Операторы сравнения}
\begin{itemize}
    \item \texttt{==} Равно
    \item \texttt{!=} Не равно
    \item \texttt{>} Больше
    \item \texttt{<} Меньше
    \item \texttt{>=} Больше или равно
    \item \texttt{<=} Меньше или равно
\end{itemize}

\subsubsection*{Логические связки}
\begin{itemize}
    \item \texttt{and} (И) — истина, только если оба условия истинны.
    \item \texttt{or} (ИЛИ) — истина, если хотя бы одно условие истинно.
    \item \texttt{not} (НЕ) — инвертирует значение (делает истину ложью и наоборот).
\end{itemize}

\begin{codefile}[python]{examples/python/04_logic.py}
\end{codefile}

\begin{taskbox}[Проверка безопасности]
Создайте переменную \texttt{wind\_speed = 7}. Напишите логическое выражение, которое возвращает \texttt{True}, если скорость ветра меньше 5 м/с.
\end{taskbox}
